% Results paragraph for the manuscript — insert in the Results section
% after describing DRN's primary classification performance.

\paragraph{Comparison with classical ML and AutoML baselines.}
To situate DRN's performance within the broader landscape of automated and
classical approaches, we benchmarked it against five classical machine
learning methods and two state-of-the-art AutoML frameworks on the same
4-class odor discrimination task (Table~\ref{tab:complexity},
Fig.~\ref{fig:complexity}).
All non-DRN methods received flattened 1,344-dimensional input vectors
(32~channels~$\times$~21~frequency bands, concatenated across OB and PCx),
as they cannot natively process the two-dimensional electrode--frequency
structure.

Among classical methods, the MLP classifier achieved the highest accuracy
(65.5\%), followed by Logistic Regression (64.5\%) and Random Forest
(63.5\%). SVM with RBF kernel reached 58.5\%, while a single Decision Tree
yielded 50.5\%---only marginally above the 25\% chance level for four
balanced classes. For AutoML, we evaluated H2O AutoML and AutoGluon, two
frameworks that consistently rank among the top performers on tabular
benchmarks. With a 5-minute training budget each, AutoGluon achieved 67.5\%
through a multi-layer stacked ensemble of LightGBM, CatBoost, and neural
network base learners, while H2O AutoML reached 54.0\% using a
gradient-boosted machine ensemble. DRN surpassed all baselines by a
substantial margin, achieving 86.5\% test accuracy---a gain of +19.0
percentage points over AutoGluon and +21.0~pp over the best classical method
(Fig.~\ref{fig:complexity}a).

This performance gap is attributable to three architectural advantages.
First, DRN's convolutional backbone preserves the spatial--spectral structure
of each brain region's band-power matrix, rather than collapsing it into a
flat vector. Second, the dual-stream design processes OB and PCx
representations independently before fusing them through a learned attention
mechanism, enabling the model to capture cross-region interactions that
single-vector approaches cannot represent. Third, DRN's inference complexity
is $O(P)$, where $P = 9.9 \times 10^6$ denotes the parameter count,
ensuring that per-sample latency (25.96~ms on CPU) remains constant
regardless of training set size---a property critical for real-time
brain--machine interface deployment. By contrast, SVM inference scales as
$O(n_{\text{sv}} \times d)$ and would degrade approximately 44-fold if the
training set grew to $n = 50{,}000$ samples, while DRN's latency would
remain unchanged (Fig.~\ref{fig:complexity}c). These results demonstrate
that domain-specific architectural inductive biases---spatial structure
preservation, dual-region fusion, and constant-time inference---yield
accuracy gains that generic AutoML pipeline search cannot recover.
